\documentclass[11pt]{article}
\usepackage[margin=1in]{geometry}
\usepackage{amsmath,amssymb,amsthm}
\usepackage[colorlinks=true,linkcolor=blue,citecolor=blue,urlcolor=blue]{hyperref}
\newtheorem{theorem}{Theorem}[section]
\newtheorem{lemma}[theorem]{Lemma}
\newtheorem{proposition}[theorem]{Proposition}
\newtheorem{corollary}[theorem]{Corollary}
\newtheorem{remark}[theorem]{Remark}
\numberwithin{equation}{section}
\newcommand{\tr}{\operatorname{tr}}
\newcommand{\rank}{\operatorname{rank}}
\newcommand{\op}{\mathrm{op}}
\newcommand{\R}{\mathbb R}
\newcommand{\C}{\mathbb C}
\newcommand{\cH}{\mathcal H}
\newcommand{\sech}{\operatorname{sech}}
\newcommand{\diag}{\operatorname{diag}}
\title{A joint endpoint law for fractional matrix powers}
\author{Working research manuscript}
\date{November 2026}
\begin{document}
\maketitle

\begin{abstract}
Let $\cH(\alpha,s)$ be the best dimension-free constant in
$\tr|B^\alpha-A^\alpha|^s\le
\cH(\alpha,s)\tr|B-A|^{\alpha s}$ for positive semidefinite matrices.
We prove the relative joint asymptotic
$\cH(1-\beta,s)\sim\beta^{2s-1}/(2\pi^2s^2)$ along every sequence
$s\downarrow0$, $\beta\downarrow0$, including arbitrarily
superexponential approaches to the endpoint. Thus
$-s\log\beta\to u$ gives the normalized crossover
$\beta s^2\cH\to e^{-2u}/(2\pi^2)$.
At a fixed Schatten exponent, the same supremum has order
$\beta^{2s-1}$ for $s<1/2$, order $\log(1/\beta)$ at $s=1/2$,
and remains bounded but strictly discontinuous at the linear endpoint
for $1/2<s<1$. A Brownian covariance kernel describes a family of
endpoint witnesses. The uniform upper uses positive spectral
compression, an explicit periodic Loewner kernel, and optimization
of the spectral masses. The nonlinear lower retains the actual
rank-one resolvent denominator. All approximation and matrix-mean
tools used in the proof have classical antecedents.
\end{abstract}

\noindent\textit{Status.}
The complete mathematical note and typeset source passed separate
independent internal reviews. External priority and external peer review
remain unresolved. No general quadratic matrix algorithm, finite-word
stability theorem, or certified novelty is asserted.

\section{Statements and context}

For $0<\alpha<1$ and $0<s\le1$, define
\begin{equation}\label{eq:H}
 \cH(\alpha,s)=
 \sup_{n\ge1}\ \sup_{\substack{A,B\in\C^{n\times n}\\
                         A,B\succeq0,\ A\ne B}}
 \frac{\tr|B^\alpha-A^\alpha|^s}{\tr|B-A|^{\alpha s}}.
\end{equation}
Write $\beta=1-\alpha$. Constants denoted $C,c>0$ are absolute unless
their dependence is indicated, and may change between occurrences.
All matrix norms without a subscript are operator norms. The symbol
$X\preceq Y$ means that $Y-X$ is positive semidefinite.

\begin{theorem}[Relative joint law]\label{thm:joint}
Along every joint sequence $s\downarrow0$ and $\beta\downarrow0$,
\begin{equation}\label{eq:joint}
 s^2\beta^{1-2s}\cH(1-\beta,s)\longrightarrow\frac1{2\pi^2}.
\end{equation}
There is no restriction on the relative rate at which the two
parameters approach zero.
\end{theorem}

\begin{corollary}[Exponential crossover]\label{cor:crossover}
If additionally $-s\log\beta\to u\in[0,\infty]$, then
\begin{equation}\label{eq:crossover}
 \beta s^2\cH(1-\beta,s)\longrightarrow
                  \frac{e^{-2u}}{2\pi^2},
 \qquad e^{-2\infty}=0.
\end{equation}
The relative asymptotic in Theorem~\ref{thm:joint} remains valid when
$u=\infty$.
\end{corollary}

For a fixed $\alpha\in(0,1)$, a companion manuscript~\cite{Companion}
proves $s^2\cH(\alpha,s)\to[2\pi^2\alpha(1-\alpha)]^{-1}$.
Formula~\eqref{eq:crossover} exhibits the failure of substituting a
moving $\alpha(s)$ in that fixed-parameter limit. At $\alpha=1$
itself the functional calculus is linear and the corresponding
constant is exactly one.

\begin{theorem}[Fixed-exponent endpoint behavior]\label{thm:fixed}
As $\beta\downarrow0$, the following statements hold.
\begin{enumerate}
\item For every fixed $0<s<1/2$ there are $c_s,C_s>0$ such that
\begin{equation}\label{eq:subcritical}
 c_s\beta^{2s-1}\le\cH(1-\beta,s)\le C_s\beta^{2s-1}
\end{equation}
for all sufficiently small $\beta$.
\item There are absolute $c,C>0$ such that
\begin{equation}\label{eq:critical}
 c\log(1/\beta)\le\cH(1-\beta,1/2)\le C\log(1/\beta)
\end{equation}
for all sufficiently small $\beta$.
\item For each fixed $1/2<s<1$, $\cH(1-\beta,s)$ stays bounded, and
\begin{equation}\label{eq:discontinuous}
 \liminf_{\beta\downarrow0}\cH(1-\beta,s)>1=\cH(1,s).
\end{equation}
\end{enumerate}
No exact fixed-$s$ multiplicative limit is asserted in
\eqref{eq:subcritical}--\eqref{eq:critical}, and
\eqref{eq:discontinuous} does not assert existence of an endpoint limit.
\end{theorem}

Dimension-free fractional-power estimates in quasi-Banach Schatten
classes are established work; see Birman--Koplienko--Solomyak,
Ricard, and Huang--Sukochev--Zanin~\cite{BKS,Ricard,HSZ}.
Ricard's Corollary~2.11 gives a dyadic divided-difference estimate
with $(1-\alpha)^{-1/p}$ dependence in its quasi-norm convention.
More directly, \cite[Remark~5.3]{Ricard} already uses the exponential
spectral scaling at $\alpha\uparrow1$, obtains the limiting
$e^{-\max(u,v)}$ kernel and its inverse-square spectral mechanism,
and identifies nonuniformity at the threshold $p=1/2$.
We do not claim those mechanisms or qualitative endpoint instability
as new. The quantitative nonlinear rates and the exact relative
joint constant are the statements examined here.
The hyperbolic-sine kernel and its positive Fourier transform below
are classical matrix-mean material~\cite{HK,Bhatia}. Positive
compression, Chebyshev approximation, Brownian covariance, the
Araki--Lieb--Thirring inequality, and Gaussian or sinc quadrature are
also classical. The specific subject here is their use to determine
the uniform nonlinear supremum in a joint singular limit. A bounded
literature comparison does not establish priority.

\section{Rank-one reduction and positive blocks}

For a positive rank-one update $E=vv^*$ with $\|v\|_2=1$, write
\[
 B=A+E,\qquad \Delta=B^\alpha-A^\alpha.
\]
Operator monotonicity gives
\begin{equation}\label{eq:contractive}
 0\preceq\Delta\preceq (A+I)^\alpha-A^\alpha\preceq I.
\end{equation}
The power triangle $\tr|X+Y|^s\le\tr|X|^s+\tr|Y|^s$, valid for
$0<s\le1$, and homogeneity imply the exact reduction
\begin{equation}\label{eq:rankone}
 \cH(\alpha,s)=\sup_{A\succeq0,\ \|v\|_2=1}
              \tr\big((A+vv^*)^\alpha-A^\alpha\big)^s.
\end{equation}
Indeed, for arbitrary positive endpoints, pass through
$A+(B-A)_+$. Telescope the positive spectral rank-one increments
from each endpoint to that common upper endpoint. Their contributions
scale by the $\alpha s$ powers of their sizes, whose sum is precisely
$\tr|B-A|^{\alpha s}$.

Put $c_\alpha=\sin(\pi\alpha)/\pi$. The fractional-power integral and
Woodbury formula give
\begin{equation}\label{eq:woodbury}
 \Delta=c_\alpha\int_0^\infty
 \frac{t^\alpha(A+tI)^{-1}vv^*(A+tI)^{-1}}{1+m(t)}\,dt,
 \qquad m(t)=v^*(A+tI)^{-1}v.
\end{equation}
Every integrand is positive. Let
\[
 P=1_{[0,1]}(A),\quad Q=I-P,\quad
 A_0=A|_{P\C^n},\quad A_1=A|_{Q\C^n},\quad v_0=Pv,\quad v_1=Qv.
\]
Compression Jensen for the concave function $t^s$ yields
\begin{equation}\label{eq:pinching}
 \tr\Delta^s\le\tr(P\Delta P)^s+\tr(Q\Delta Q)^s.
\end{equation}
In the first compression, the denominator in~\eqref{eq:woodbury}
is at least $1+v_0^*(A_0+tI)^{-1}v_0$. Therefore
\begin{equation}\label{eq:lowdom}
 P\Delta P\preceq
 \Delta_0:=(A_0+v_0v_0^*)^\alpha-A_0^\alpha,
 \qquad 0\preceq A_0\preceq I,\quad\|v_0\|_2\le1.
\end{equation}
For the high compression, deleting the denominator gives
\begin{equation}\label{eq:highdom}
 Q\Delta Q\preceq L:=D(x^\alpha)_{A_1}[v_1v_1^*],\qquad A_1\succeq I.
\end{equation}
Here $D$ denotes the Fr\'echet derivative. In a spectral basis
$A_1=\diag(\lambda_i)$,
\begin{equation}\label{eq:loewnerentries}
 L_{ij}=(v_1)_i\overline{(v_1)_j}
 \frac{\lambda_i^\alpha-\lambda_j^\alpha}{\lambda_i-\lambda_j},
\end{equation}
with the diagonal quotient interpreted as $\alpha\lambda_i^{\alpha-1}$.

\section{Uniform quadrature estimates}

We give the uniformity needed when a decay slope tends to zero.
The strip trapezoidal inequality used below is
\begin{equation}\label{eq:trapezoidal}
 \left\|\int_\R G(x)\,dx-h\sum_{n\in\mathbb Z}G(nh)\right\|
 \le\frac{2M}{e^{2\pi d/h}-1}
\end{equation}
if both boundary $L_1$ norms on $\Im z=\pm d$ are at most $M$.
The scalar contour proof followed by duality proves this for
Banach-valued functions; see also~\cite{TW}.

\begin{lemma}[Uniform tapered rule]\label{lem:taper}
Fix $d>0$. Suppose $0<\beta\le1/2$, $\alpha=1-\beta$, and $F$ is
Banach-valued, analytic on a neighborhood of $|\Im z|\le d$, with
\begin{equation}\label{eq:smallenvelope}
 \|F(x+i\eta)\|\le M\beta e^{-g(x)},\qquad
 g(x)=\alpha(-x)_++\beta x_+,
 \qquad |\eta|\le d.
\end{equation}
There are positive-weight real-node quadratures of error $e^{-L}$,
for $L\ge L_{M,d}$ independent of $\beta$, with node count
\begin{equation}\label{eq:taperedcount}
 q_L\le\frac{L^2}{4\pi d\alpha\beta}
                   +C_{M,d}\beta^{-1}(L+1)^{7/4}.
\end{equation}
If each node value has rank at most $r$, the quadrature has rank at
most $rq_L$.
\end{lemma}
\begin{proof}
Put $w=L^{3/4}$, $I_j=[jw,(j+1)w]$, and
\[
 \phi_j(z)=\frac1{\sqrt\pi}\int_{I_j}e^{-(z-u)^2}\,du,
 \qquad m_j=\min_{I_j}g.
\]
The windows are entire, positive on the real line, and sum to one.
Since $g$ is $1$-Lipschitz and
$|\phi_j(x+i\eta)|\le e^{d^2}\phi_j(x)$, the two boundary
integrals of $F\phi_j$ are at most $C_{M,d}\beta w e^{-m_j}$.
Retain the cells meeting
$[-(L+w)/\alpha,(L+w)/\beta]$; their number is
$J\le C L/(\beta w)+C$. The discarded-window integral is bounded by
\[
 C_{M,d}\beta(\alpha^{-1}+\beta^{-1})e^{-L-w}
 \le C_{M,d}e^{-L-w}.
\]
This follows by viewing the omitted windows as convolution with the
Gaussian and using $g(x)\ge g(u)-|x-u|$.

For each retained cell use mesh
\[
 Q_j=\max(w,L-m_j+3w),\qquad h_j=2\pi d/Q_j.
\]
Equation~\eqref{eq:trapezoidal} and $m_j+Q_j\ge L+3w$ give total
infinite-mesh error at most
$C_{M,d}\beta wJ e^{-L-3w}\le C_{M,d}L e^{-L-3w}$.
Keep only nodes in $I_j+[-2\sqrt L,2\sqrt L]$.
The Gaussian bound on $\phi_j$ and $h_j\le2\pi d/w$ give discarded
lattice mass at most $C_{M,d}\beta w e^{-2L}$ per cell.
Its sum is at most $C_{M,d}L e^{-2L}$. Thus the total error is
at most $e^{-L}$ beyond a threshold independent of $\beta$.

The node count is bounded by
\[
 \sum_j\left(\frac{(w+4\sqrt L)Q_j}{2\pi d}+2\right).
\]
The Riemann sum for $(L-g)_+$ differs from its integral by
$O(\beta^{-1}Lw)$, and
\[
 \int_\R(L-g(x))_+\,dx=\frac{L^2}{2}
                     (\alpha^{-1}+\beta^{-1})=\frac{L^2}{2\alpha\beta}.
\]
The $3w$ overhead is $O(\beta^{-1}Lw)$; enlarging each cell adds
$O(\beta^{-1}L^{5/2}/w)$ nodes, and rounding adds $O(J)$.
With $w=L^{3/4}$ this proves~\eqref{eq:taperedcount}. Every weight is
$h_j\phi_j(nh_j)>0$. Rank subadditivity proves the last claim.
\end{proof}

For positive $A,B$ with $\|B-A\|\le1$, their fractional difference
also has the resolvent form
\begin{equation}\label{eq:resdiff}
 B^\alpha-A^\alpha=c_\alpha\int_0^\infty
              t^\alpha[(A+tI)^{-1}-(B+tI)^{-1}]\,dt.
\end{equation}
For fixed $d<\pi$ and $t=\rho e^{i\eta}$, $|\eta|\le d$,
$|a+t|\ge\cos(d/2)(a+\rho)$ for $a\ge0$. Bounding the resolvents
separately and then by their difference identity gives
\[
 \|(A+tI)^{-1}-(B+tI)^{-1}\|
              \le C_d\min(\rho^{-1},\rho^{-2}).
\]
Since $c_\alpha=\sin(\pi\beta)/\pi\le\beta$, the logarithmic
integrand in~\eqref{eq:resdiff} has envelope~\eqref{eq:smallenvelope}
with $M=C_d$ independent of $\beta$.
For rank-one pairs, Lemma~\ref{lem:taper} and best-rank approximation
therefore imply, writing $N_\Delta(L)$ for the number of singular
values exceeding $e^{-L}$,
\begin{equation}\label{eq:uniformcount}
 N_\Delta(L)\le\frac{L^2}{4\pi d\alpha\beta}
                    +C_d\beta^{-1}(L+1)^{7/4}
\end{equation}
beyond a uniform threshold. A fixed enlargement bounds the finite
initial range by $C_d\beta^{-1}$. Because $\|\Delta\|\le1$, layer cake
gives
\begin{equation}\label{eq:gaussianupper}
 \beta s^2\cH(1-\beta,s)
       \le\frac1{2\pi d\alpha}+C_d s^{1/4},
 \qquad 0<s\le1/4,\quad0<\beta\le1/2.
\end{equation}

We also need a coarser count valid with an absolute constant:
\begin{equation}\label{eq:coarsecount}
 N_\Delta(L)\le C\beta^{-1}(L+1)^2,\qquad L\ge0.
\end{equation}
For completeness, choose a fixed strip $d<\pi$, mesh
$h=c/(L+C_0)$, and truncate the sinc rule between logarithmic
endpoints $-(L+C_0)/\alpha$ and $(L+C_0)/\beta$.
The boundary $L_1$ norms are
$O(c_\alpha(\alpha^{-1}+\beta^{-1}))=O(1)$.
The two lattice tails are $O(e^{-L-C_0})$, including the
$h/(1-e^{-\beta h})$ factor because $c_\alpha/\beta\le1$.
Equation~\eqref{eq:trapezoidal} proves the error bound and
$O(\beta^{-1}(L+1)^2)$ nodes prove~\eqref{eq:coarsecount}.
In particular,
\begin{equation}\label{eq:coarseeigen}
 \lambda_j(\Delta)\le C e^{-c\sqrt{\beta j}}.
\end{equation}

\begin{lemma}[Uniform low block]\label{lem:low}
For $1/2\le\alpha<1$, the corrections $\Delta_0$ in
\eqref{eq:lowdom} satisfy
\begin{equation}\label{eq:lowmoment}
 \tr\Delta_0^s\le C/s^2,\qquad0<s\le1,
\end{equation}
with $C$ independent of $\alpha$, dimension, and the endpoints.
\end{lemma}
\begin{proof}
Write $E_0=v_0v_0^*$. Subtract the rank-one integral
\[
 c_\alpha\int_0^\infty\frac{t^\alpha}{(1+t)^2}\,dt\ E_0
 =\alpha E_0.
\]
The residual logarithmic integrand is
\[
 F(z)=c_\alpha t^{\alpha+1}
 \big[(A_0+tI)^{-1}-(A_0+E_0+tI)^{-1}-E_0/(1+t)^2\big],\quad t=e^z.
\]
For a fixed $d<\pi$, it obeys the common bound
\[
 \|F(x+i\eta)\|\le C_d
 \begin{cases}e^{x/2},&x\le0,\\ e^{-x},&x\ge0,\end{cases}
 \qquad |\eta|\le d.
\]
Indeed small $|t|$ uses separate resolvent bounds. For $|t|\ge1$,
both endpoints have norm at most two, and
$(A_0+tI)^{-1}E_0(A_0+E_0+tI)^{-1}=E_0/t^2+O_d(|t|^{-3})$;
the subtracted rational term has the same leading term.
Every residual node has rank at most two. A fixed-mesh sinc rule
with the displayed fixed slopes gives error $e^{-L}$ using
$C(L+1)^2$ nodes, uniformly in $\alpha$. Including $\alpha E_0$
changes the rank by at most one. Since $\|\Delta_0\|\le1$, layer
cake proves~\eqref{eq:lowmoment}.
\end{proof}

\section{The periodic Loewner kernel and its corners}

For the high block put $x_i=\log\lambda_i\ge0$ and $z_i=(v_1)_i$.
Then~\eqref{eq:loewnerentries} is
\begin{equation}\label{eq:kernel}
 L_{ij}=z_i\overline{z_j}e^{-\beta(x_i+x_j)/2}k_\beta(x_i-x_j),
 \qquad k_\beta(x)=\frac{\sinh((1-\beta)x/2)}{\sinh(x/2)}.
\end{equation}
The value at zero is $1-\beta$. The classical transform identity is
\begin{equation}\label{eq:transform}
 \widehat k_\beta(\xi):=\int_\R k_\beta(x)e^{-i\xi x}\,dx
       =\frac{2\pi\sin(\pi\beta)}{\cosh(2\pi\xi)-\cos(\pi\beta)}.
\end{equation}
For example, expanding on the two half-lines and integrating gives
the partial-fraction expansion of the cotangent and hence this
identity. It is also the Poisson kernel for a strip after rescaling.
Its strict positivity makes $k_\beta$ positive definite.

For $\Re z\ge0$ separate its slowly decaying part:
\begin{equation}\label{eq:remainder}
 k_\beta(z)=e^{-\beta z/2}+r_\beta(z),\qquad
 r_\beta(z)=-e^{-(1-\beta/2)z}
                      \frac{1-e^{-\beta z}}{1-e^{-z}}.
\end{equation}
The value at zero is removable. Uniformly for $0<\beta\le1/2$,
$\Re z\ge-1/4$, and $|\Im z|\le1/4$,
\begin{equation}\label{eq:rbound}
 |r_\beta(z)|\le C\beta(1+(\Re z)_+)e^{-3(\Re z)_+/4}.
\end{equation}
To see this, write
$(1-e^{-\beta z})/\beta=z\int_0^1e^{-t\beta z}\,dt$.
On a bounded part of this half-strip $(1-e^{-z})/z$ is bounded away
from zero; on the remaining part $|1-e^{-z}|$ is bounded away from
zero. The exponent $1-\beta/2$ is at least $3/4$.

\begin{lemma}[Uniform corner estimate]\label{lem:corner}
Suppose $h$ is holomorphic on a neighborhood of
$\Re z\ge-1/4$, $|\Im z|\le1/4$, and
\[
 |h(z)|\le C\beta(1+(\Re z)_+)e^{-c(\Re z)_+}.
\]
For arbitrary finite nonnegative nodes $a_i,b_j$ and vectors $p,q$,
let $T_{ij}=p_i h(a_i+b_j)\overline{q_j}$. Then
\begin{equation}\label{eq:cornermoment}
 \tr|T|^s\le\frac{C'}{s^2}
                 \big(\beta\|p\|_2\|q\|_2\big)^s,\qquad0<s\le1.
\end{equation}
The constant is independent of nodes, dimensions, $\beta$, and
the lengths of the node intervals.
\end{lemma}
\begin{proof}
Fix $L\ge1$ and $R=C_0(L+1)$. Outside the square $a,b\le R$ the
kernel has magnitude at most $\beta e^{-L}$, if $C_0$ is sufficiently
large. On $[0,2R]$, a Bernstein ellipse with parameter
$\rho=1+c_0/R$, after affine rescaling, stays in the half-strip in
the statement. Its supremum is at most $C_1\beta$.
The Chebyshev error of degree $N$ is at most
$C_2\beta R e^{-c_1N/R}$. Choose
$N=C_3R(L+\log(R+1))\le C_4(L+1)^2$ to make it at most
$\beta e^{-L}$.

Restrict the resulting polynomial kernel $p_N(a+b)$ to the square
and set it to zero outside. Collecting powers of $a$ writes it as
$\sum_{j=0}^N a^j q_j(b)$, so its rank is at most $N+1$.
The weighted residual has Frobenius norm at most
$\beta e^{-L}\|p\|_2\|q\|_2$. The original operator norm is at most
$C\beta\|p\|_2\|q\|_2$. Counting singular values and integrating
layer cake gives~\eqref{eq:cornermoment}.
\end{proof}

If $\ell\ge1$, the function
\begin{equation}\label{eq:himage}
 h_{\beta,\ell}(z)=\sum_{n\ge0}r_\beta(z+n\ell)
\end{equation}
satisfies the same bound as in Lemma~\ref{lem:corner}, with absolute
constants independent of $\beta,\ell$. This follows from
\eqref{eq:rbound} and uniform geometric summation, including local
uniform convergence of the analytic series.

\begin{lemma}[One spectral cell]\label{lem:cell}
Let the nodes $x_i$ in~\eqref{eq:kernel} lie in a single cell of
length $\ell\ge1$. Put $m=\sum_i|z_i|^2e^{-\beta x_i}$, and assume
$\beta\ell\le1$. For every $0<s<1/2$,
\begin{equation}\label{eq:cellbound}
 \tr L^s\le m^s\left[
 \beta^s\ell^{1-s}\frac{(\pi^2)^s I_s}{\pi^2}
       +C_1\left(\frac{C_2}{\beta\ell}\right)^s
       +\frac{C_3\beta^s}{s^2}\right],
 \qquad I_s=2^{2s-1}B(s,1-2s).
\end{equation}
The constants $C_i$ are absolute.
\end{lemma}
\begin{proof}
In local coordinates $y_i\in[0,\ell)$, periodize the kernel:
\[
 p_{\beta,\ell}(z)=\sum_{n\in\mathbb Z}k_\beta(z+n\ell)
       =\sum_{k\in\mathbb Z}d_k e^{2\pi ikz/\ell},
 \quad
 d_k=\frac{2\pi\sin(\pi\beta)}
 {\ell[\cosh(4\pi^2k/\ell)-\cos(\pi\beta)]}>0.
\]
The Fourier series converges absolutely. Write
$w_i=z_i e^{-\beta x_i/2}$ and replace $k_\beta$ by this
periodization to obtain a positive matrix $P$. Each Fourier term is
positive rank one, with its generating vector of squared norm $m$.
The power triangle therefore gives
\begin{equation}\label{eq:periodicpower}
 \tr P^s\le m^s\sum_{k\in\mathbb Z}d_k^s.
\end{equation}
No orthogonality of Fourier vectors in the arbitrary spectral
measure is assumed. The zero mode satisfies
\[
 d_0=2\pi\cot(\pi\beta/2)/\ell\le4/(\beta\ell).
\]
For $k\ne0$, use $\sin(\pi\beta)\le\pi\beta$ and
$\cosh q-\cos(\pi\beta)\ge2\sinh^2(q/2)$ to obtain
\[
 d_k\le\frac{\pi^2\beta}{\ell}
                  \operatorname{csch}^2(2\pi^2k/\ell).
\]
The summand is decreasing, and its integral from zero is finite for
$s<1/2$. Consequently
\begin{align}\label{eq:fouriersum}
 \sum_{k\ne0}d_k^s
 &\le 2\left(\frac{\pi^2\beta}{\ell}\right)^s
       \int_0^\infty\sinh(2\pi^2x/\ell)^{-2s}\,dx\nonumber\\
 &=\beta^s\ell^{1-s}\frac{(\pi^2)^s}{\pi^2}
                   \int_0^\infty(\sinh u)^{-2s}\,du.
\end{align}
Substitution $t=e^{-2u}$ evaluates the last integral as
$I_s=2^{2s-1}B(s,1-2s)$.

It remains to charge the images removed from the periodic kernel.
For $|z|\le\ell$, their slowly decaying parts in~\eqref{eq:remainder}
sum to
\[
 \frac{e^{-\beta\ell/2}}{1-e^{-\beta\ell/2}}
                         (e^{-\beta z/2}+e^{\beta z/2}).
\]
After weighting, these form a matrix $Z$ of rank at most two and
norm at most $Cm/(\beta\ell)$. Indeed the coefficient is at most
$2/(\beta\ell)$ and the two local multipliers
$e^{\pm\beta y_i/2}$ have norms bounded by $e^{1/2}$.
The remaining images are two Hankel corners
$h_{\beta,\ell}(\ell+y_i-y_j)$ and
$h_{\beta,\ell}(\ell-y_i+y_j)$. In each, reflect one coordinate so
the argument is a sum of two nonnegative coordinates.
Lemma~\ref{lem:corner} bounds their combined power trace by
$C\beta^sm^s/s^2$. Reflection relabels the nodes and preserves
the Euclidean norms of the two weighted vectors.

The exact identity is $L=P-Z-H_1-H_2$. Apply the power triangle to
these possibly indefinite image corrections, use the rank bound on
$Z$, and combine~\eqref{eq:periodicpower}--\eqref{eq:fouriersum}.
This proves~\eqref{eq:cellbound}. The whole image tail, not only
its nearest term, was included in~\eqref{eq:himage}.
\end{proof}

\section{Mass optimization and the uniform small-endpoint upper}

Partition $x\ge0$ into $I_j=[j\ell,(j+1)\ell)$. Write
\[
 \mu_j=\sum_{x_i\in I_j}|z_i|^2,\qquad
 m_j=\sum_{x_i\in I_j}|z_i|^2e^{-\beta x_i}
                  \le e^{-\beta j\ell}\mu_j.
\]
We have $\sum_j\mu_j\le1$. Positive pinching gives
$\tr L^s\le\sum_j\tr L_j^s$ over the finitely many occupied cells.
Ordinary H\"older with exponents $1/s$ and $1/(1-s)$ gives
\begin{equation}\label{eq:massopt}
 \sum_jm_j^s\le\sum_j\mu_j^se^{-s\beta j\ell}
 \le\left(\sum_{j\ge0}e^{-s\beta j\ell/(1-s)}\right)^{1-s}
 =[1-e^{-s\beta\ell/(1-s)}]^{s-1}.
\end{equation}
This step is uniform over all atomic measures and dimensions.

Take $\ell=1/\beta$, so the one-cell hypothesis
$\beta\ell\le1$ holds with equality. Define
\begin{equation}\label{eq:AsTs}
 T_s=[1-e^{-s/(1-s)}]^{s-1},\qquad
 A_s=\frac{(\pi^2)^s I_s}{\pi^2}.
\end{equation}
Their elementary limits are
\begin{equation}\label{eq:TsIlimits}
 sT_s\longrightarrow1,\qquad sI_s\longrightarrow\frac12,
 \qquad sA_s\longrightarrow\frac1{2\pi^2}.
\end{equation}
Combining~\eqref{eq:pinching}, Lemmas~\ref{lem:low} and
\ref{lem:cell}, and~\eqref{eq:massopt} proves the nonasymptotic estimate
\begin{equation}\label{eq:uniformupper}
 \cH(1-\beta,s)\le\frac C{s^2}
       +T_s\left[A_s\beta^{2s-1}+C+\frac{C\beta^s}{s^2}\right],
 \qquad0<\beta\le1/2,\quad0<s<1/2.
\end{equation}
All constants outside $A_s,T_s$ are absolute. This immediately
gives the upper half of~\eqref{eq:subcritical} for fixed $s<1/2$.
For the joint limit multiply by $s^2\beta^{1-2s}$:
\begin{align}\label{eq:relativeupper}
 s^2\beta^{1-2s}\cH(1-\beta,s)
 \le {}&s^2T_sA_s+C\beta^{1-2s}
       +Cs^2T_s\beta^{1-2s}+CT_s\beta^{1-s}.
\end{align}
If $\beta\le s^3$, $0<s\le1/4$, the three errors are respectively
$O(s^{3-6s})$, $O(s^{4-6s})$, and $O(s^{2-3s})$.
They tend to zero uniformly even when $\beta$ is superexponentially
small. This explicitly controls the relative normalization.

\section{A matching nonlinear finite-matrix lower}

The following convolution inequality is useful for retaining the
nonlinear denominator. Let $I$ be a bounded interval,
\[
 C(x-y)=\frac1{2\cosh((x-y)/2)},\qquad
 h(\xi)=\widehat C(\xi)=\pi\sech(\pi\xi),\qquad
 C_I=1_I C1_I.
\]
If $a,b$ are continuous and bounded above and below by positive
numbers on $I$, then, for $0<p\le1$,
\begin{equation}\label{eq:convlower}
 \tr|M_aC_IM_b|^p\ge\frac1{2\pi}
           \int_I(a(x)b(x))^p\,dx\int_\R h(\xi)^p\,d\xi.
\end{equation}
Indeed $M_aC_IM_b$ is similar to the positive operator
$S=M_{\sqrt{ab}}C_IM_{\sqrt{ab}}$. Multiplicative eigenvalue--singular-value
majorization gives $\tr|M_aC_IM_b|^p\ge\tr S^p$.
The reversed Araki--Lieb--Thirring inequality~\cite{ALT} and
compression Jensen~\cite{HP} give
\[
 \tr S^p\ge\tr(M_{(ab)^{p/2}}C_I^pM_{(ab)^{p/2}}),\qquad
 C_I^p\succeq1_I C^p1_I.
\]
Taking the weighted trace and using Plancherel proves
\eqref{eq:convlower}. The compact analytic kernel on $I$ belongs
to every Schatten class; the intermediate positive traces also
follow by finite-rank approximation or monotone extended traces.

Fix $\eta=1/2$, $\delta>0$, and the probability density
\begin{equation}\label{eq:density}
 \rho(x)=\begin{cases}c_0e^{\eta x},&x\le0,\\
                      c_0e^{-\delta x},&x\ge0,
        \end{cases}\qquad
 c_0=(\eta^{-1}+\delta^{-1})^{-1}.
\end{equation}
Put $m(t)=\int_\R\rho(x)/(e^x+t)\,dx$. On $L_2(I)$ let $D_I$
multiply by $e^x$ and let $w_I=\sqrt\rho|_I$. The input increment
$w_Iw_I^*$ has norm at most one. For
$\Delta_I=(D_I+w_Iw_I^*)^\alpha-D_I^\alpha$, Woodbury gives the
integral~\eqref{eq:woodbury}, with $m_I(t)\le m(t)$.
Replacing $m_I$ by $m$ lowers that positive integral, as does
restricting $t=e^y$ to $y\in I$. Consequently
\begin{equation}\label{eq:lowerfactor}
 \Delta_I\succeq T_IT_I^*,\qquad
 T_I=M_aC_IM_b,
\end{equation}
where
\[
 a(x)=\sqrt{c_\alpha\rho(x)}e^{-x/2},\quad
 b(x)=\frac{e^{\alpha x/2}}{\sqrt{1+m(e^x)}},\quad
 q(x)=(a(x)b(x))^2
       =\frac{c_\alpha\rho(x)e^{-\beta x}}{1+m(e^x)}.
\]
Direct substitution using
$C(x-y)=e^{(x+y)/2}/(e^x+e^y)$ verifies
\eqref{eq:lowerfactor}; no false global compression ordering is used.

We explain explicitly why these bounded-interval operators give a
lower bound for the finite-matrix supremum~\eqref{eq:H}.
Approximate $e^x$ uniformly on $I$ by positive step functions $D_n$,
keeping $w_I$ fixed. Continuous functional calculus gives convergence
of the corrections in operator norm. The span of the finitely many
cell vectors $1_Jw_I$ reduces both endpoints; outside it the two
endpoints agree. Thus every approximating correction is a finite
matrix correction extended by zero. Its input norm is
$t_I=\|w_I\|_2^2\le1$, so its normalized ratio is at least its
output power trace. The limit is compact by norm approximation by
finite-rank corrections. Every finite partial sum of singular-value
$s$-powers passes to the limit, and then their supremum does.
Thus $\cH(\alpha,s)\ge\tr\Delta_I^s$.

Applying~\eqref{eq:convlower} at $p=2s\le1$, then exhausting the
line only in the resulting scalar integral, proves
\begin{equation}\label{eq:nonlinearlower}
 \cH(\alpha,s)\ge\frac{J_s}{2\pi}\int_\R q(x)^s\,dx,
 \qquad
 J_s=\int_\R[\pi\sech(\pi\xi)]^{2s}\,d\xi.
\end{equation}
In particular, the actual nonlinear denominator has been retained
before the finite-matrix transfer. A beta integral yields
\begin{equation}\label{eq:J}
 J_s=\pi^{2s-1}\frac{\sqrt\pi\,\Gamma(s)}{\Gamma(s+1/2)},
 \qquad sJ_s\longrightarrow\frac1\pi.
\end{equation}

Choose now $\delta=s\beta/(1-s)$. For $x\ge0$,
$m(e^x)\le e^{-x}\le1$, hence
\[
 \int_\R q(x)^s\,dx
       \ge\frac{(c_\alpha c_0/2)^s}{s(\beta+\delta)}.
\]
For $s\le1/4$, $\beta\le1/2$, we have $\delta\le1/6$ and
\[
 \frac{c_\alpha c_0}{2}\ge
          \beta^2\frac{3s}{4\pi(1-s)}.
\]
It follows that
\begin{equation}\label{eq:lowerbs}
 \beta s^2\cH(1-\beta,s)\ge b_s\beta^{2s},\qquad
 b_s=\frac{sJ_s}{2\pi}(1-s)
                  \left(\frac{3s}{4\pi(1-s)}\right)^s
       \longrightarrow\frac1{2\pi^2}.
\end{equation}
This lower bound is uniform over all $\beta\in(0,1/2]$.

\begin{proof}[Proof of Theorem~\ref{thm:joint}]
For $\beta\le s^3$, the upper~\eqref{eq:relativeupper} and
lower~\eqref{eq:lowerbs}, after dividing by $\beta^{2s}$,
give the desired relative normalized limit uniformly.
For $\beta\ge s^3$ along a joint sequence with $\beta\to0$,
$\beta^{2s}\to1$. The lower tends to $1/(2\pi^2)$.
For any fixed $d<\pi$,~\eqref{eq:gaussianupper} gives the upper
limit $1/(2\pi d)$, since $\alpha\to1$. Taking $d\uparrow\pi$
proves the matching upper. Every sequence splits into these two
regimes. Corollary~\ref{cor:crossover} follows immediately.
\end{proof}

\section{Brownian covariance at the endpoint}

The scaling and the equivalent limiting kernel in this section occur
in \cite[Remark~5.3]{Ricard}. We record their nonlinear finite-matrix
realization and the uniform quantitative comparisons needed below.

\begin{proposition}[Nonlinear endpoint witnesses]\label{prop:brownianlimit}
Fix finitely many $u_i>0$ and a unit vector $v$. Set
$A_\beta=\diag(e^{u_i/\beta})$, $B_\beta=A_\beta+vv^*$.
Then in operator norm
\begin{equation}\label{eq:brownianlimit}
 B_\beta^{1-\beta}-A_\beta^{1-\beta}\longrightarrow K,
 \qquad K_{ij}=v_i\overline{v_j}\min(t_i,t_j),\quad t_i=e^{-u_i}.
\end{equation}
\end{proposition}
\begin{proof}
Let $L_\beta=D(x^{1-\beta})_{A_\beta}[vv^*]$ and
$a_{\min}=\min_i e^{u_i/\beta}$. Since $m(t)\le a_{\min}^{-1}$,
\eqref{eq:woodbury} gives
\[
 (1+a_{\min}^{-1})^{-1}L_\beta
 \preceq B_\beta^{1-\beta}-A_\beta^{1-\beta}\preceq L_\beta.
\]
Also $\|L_\beta\|\le(1-\beta)a_{\min}^{-\beta}\le1$.
For unequal $u_i,u_j$, the divided difference tends to
$e^{-\max(u_i,u_j)}$ by factoring out the larger exponential.
On the diagonal it is $(1-\beta)e^{-u_i}$, with the same limit.
Finite entrywise convergence and $a_{\min}^{-1}\to0$ prove the claim.
\end{proof}

We next give a uniform upper counterpart for arbitrary high-base
spectral measures. With $u_i=\beta\log\lambda_i$,
\eqref{eq:kernel} becomes a weighted kernel
$k_\beta((u_i-u_j)/\beta)$. Its Fourier transform is
\begin{equation}\label{eq:scaledtransform}
 \frac{2\pi\beta\sin(\pi\beta)}
              {\cosh(2\pi\beta\xi)-\cos(\pi\beta)}
       \le\frac1{\xi^2+1/4}.
\end{equation}
Indeed $\cosh(2\pi\beta\xi)-1\ge2\pi^2\beta^2\xi^2$
and $\sin(\pi\beta)\le\pi\beta$, while
$1-\cos(\pi\beta)\ge(\pi\beta/2)\sin(\pi\beta)$ follows from
$\tan(x/2)\ge x/2$. The right side is the transform of
$e^{-|u|/2}$. Fourier positivity therefore proves
\begin{equation}\label{eq:browniandom}
 0\preceq L\preceq K,\qquad
 K_{ij}=z_i\overline{z_j}\min(t_i,t_j),\quad t_i=e^{-u_i}\in(0,1],
 \quad\sum_i|z_i|^2\le1.
\end{equation}

\begin{lemma}[Uniform Brownian eigenvalues]\label{lem:browneigen}
Every matrix $K$ in~\eqref{eq:browniandom} satisfies
\begin{equation}\label{eq:browneigen}
 \lambda_j(K)\le C/j^2
\end{equation}
with an absolute constant.
\end{lemma}
\begin{proof}
Let $\mu=\sum_i|z_i|^2\delta_{t_i}$, combining equal atoms if
necessary, and define $T:L_2(\mu)\to L_2([0,1],dx)$ by
$Tf(x)=\int_{[x,1]}f(t)\,d\mu(t)$. Its $T^*T$ is the Brownian
Gram operator, with the nonzero spectrum of $K$.
Partition $[0,1]$ at the quantiles with mass levels $2k/n$ of
Lebesgue measure plus $\mu$, including the endpoints and deduplicating
repeated cuts at atoms. Every open interval between successive cuts
has both length and interior $\mu$ mass at most $2/n$.
On $(a,b)$ approximate $Tf$ by the constant $\int_{[b,1]}f\,d\mu$.
Boundary atoms are thus included exactly. This step-function map has
rank at most $n+1$. The residual is a direct sum of local Volterra
maps, each of norm at most
$\sqrt{|b-a|\mu((a,b))}\le2/n$.
Thus the $(n+2)$nd singular value of $T$ is at most $2/n$, and
$\lambda_{n+2}(K)\le4/n^2$. Enlarging the constant covers the first
few eigenvalues, using $\tr K\le1$.
\end{proof}

For every fixed $s>1/2$,~\eqref{eq:pinching},
Lemma~\ref{lem:low}, and~\eqref{eq:browneigen} give a bound on
$\cH(1-\beta,s)$ independent of $\beta$.
To see strict discontinuity for $s<1$, first take the Brownian
two-node model with $t=(1,1-d)$, $v=(1,1)/\sqrt2$ and $0<d\le1/2$.
Its larger eigenvalue is at least $1-d$, and its smaller eigenvalue
is at least $d/8$: the determinant is $d(1-d)/4$ and the larger
eigenvalue is at most one. Hence its power trace is at least
\[
 (1-d)^s+(d/8)^s\ge1-d+(d/8)^s>1
\]
for sufficiently small $d$ depending on $s<1$.
A common scaling $t\mapsto\tau t$, with $\tau<1$ sufficiently
close to one, preserves the strict inequality and makes every $u_i$
strictly positive. Proposition~\ref{prop:brownianlimit} now proves
\eqref{eq:discontinuous}.

\section{The critical exponent and matching fixed-exponent orders}

We give a quantitative family of Brownian witnesses. Fix
$0<a<b<1$, let $\mu$ be uniform probability on $t\in[a,b]$, and
write $u=-\log t$. The derivative $L_\beta$ and Brownian operator
$K$ have kernels as in the preceding section on $L_2(\mu)$.
Their difference is positive. Its convolution factor is
\begin{equation}\label{eq:diffkernel}
 d_\beta(h)=e^{-|h|/2}-k_\beta(h/\beta)
 =e^{-|h|/2}\frac{e^{|h|}-1}{e^{|h|/\beta}-1}\ge0.
\end{equation}
For $\beta\le1/2$,
\begin{equation}\label{eq:dL1}
 \int_\R d_\beta(h)\,dh\le\pi^2\beta^2.
\end{equation}
In fact, on $h\ge0$ use $e^h-1\le he^h$ and put $h=\beta z$.
The resulting integral is bounded by
$\beta^2\int_0^\infty z e^{-z/2}/(1-e^{-z})\,dz
=\pi^2\beta^2/2$; double it for the other half-line.
The $u$ density of $\mu$ is at most $b/(b-a)$ and the exterior
weight $e^{-(u+v)/2}$ is at most $b$. Schur's test gives
\begin{equation}\label{eq:diffnorm}
 \|K-L_\beta\|\le\frac{\pi^2b^2}{b-a}\beta^2.
\end{equation}

The Brownian kernel $\min(t,t')$ equals the positive rank-one
constant $a$ plus the Brownian kernel translated to $[0,b-a]$.
On uniform probability on $[0,\ell_0]$, that latter kernel has
eigenvalues
\begin{equation}\label{eq:explicitbrown}
 \frac{\ell_0}{\pi^2(j-1/2)^2},\qquad j\ge1.
\end{equation}
For example, differentiating the integral eigenvalue equation twice
gives a Dirichlet condition at zero and a Neumann condition at
$\ell_0$. By min--max,~\eqref{eq:diffnorm} implies
\begin{equation}\label{eq:manybrown}
 \lambda_j(L_\beta)\ge c/j^2,
 \qquad 1\le j\le c'/\beta,
\end{equation}
where $c,c'>0$ depend only on $a,b$.

For the actual nonlinear correction let the base multiply by
$e^{u/\beta}$ and let the unit perturbation vector be the constant
one in $L_2(\mu)$. This is a bounded positive base for every fixed
$\beta$. As in Proposition~\ref{prop:brownianlimit}, the correction
dominates $(1+b^{1/\beta})^{-1}L_\beta$.
All lower bounds transfer to finite matrices by positive step
approximation of this bounded multiplication operator and convergence
of finite partial singular-value sums, exactly as in the earlier
nonlinear lower transfer. Summing~\eqref{eq:manybrown} proves
$\cH(1-\beta,s)\ge c_s\beta^{2s-1}$ for every fixed $s<1/2$,
and $\cH(1-\beta,1/2)\ge c\log(1/\beta)$.

For the critical upper, the high compression satisfies both
\eqref{eq:coarseeigen} and~\eqref{eq:browneigen}, by min--max and
\eqref{eq:highdom}. Split its square-root trace at
\[
 J=\left\lceil C_0\beta^{-1}\log^2(e/\beta)\right\rceil.
\]
The first part is at most $C\sum_{j\le J}j^{-1}
\le C\log(e/\beta)$. For sufficiently large fixed $C_0$, the tail
from~\eqref{eq:coarseeigen} is bounded by
\[
 C\int_J^\infty e^{-c\sqrt{\beta x}/2}\,dx
 \le C\beta^{-1}(\sqrt{\beta J}+1)e^{-c\sqrt{\beta J}/2}
 \le C.
\]
The low block contributes an absolute constant by
Lemma~\ref{lem:low}. This proves~\eqref{eq:critical}.
Together with the upper~\eqref{eq:uniformupper} for fixed $s<1/2$
and the bounded/discontinuous argument above, it completes
Theorem~\ref{thm:fixed}.

\begin{thebibliography}{99}
\bibitem{BKS}
M. Sh. Birman, L. S. Koplienko, M. Z. Solomyak.
Estimates of the spectrum of a difference of fractional powers of
selfadjoint operators. Izv. Vyssh. Uchebn. Zaved. Mat., no.~3,
3--10, 1975. \url{https://www.mathnet.ru/eng/ivm6330}.

\bibitem{Ricard}
\'E. Ricard. Fractional powers on noncommutative $L_p$ for $p<1$.
Adv. Math. 333, 194--211, 2018.
\url{https://arxiv.org/abs/1805.05677}.

\bibitem{HSZ}
J. Huang, F. Sukochev, D. Zanin. Operator $\theta$-H\"older functions
with respect to $\|\cdot\|_p$, $0<p\le\infty$.
J. London Math. Soc., 2022. \url{https://arxiv.org/abs/2110.09708}.

\bibitem{HK}
F. Hiai, H. Kosaki. Means of Hilbert Space Operators.
Lecture Notes in Mathematics 1820, Springer, 2003.

\bibitem{Bhatia}
R. Bhatia. Positive Definite Matrices. Princeton University Press, 2007.

\bibitem{TW}
L. N. Trefethen, J. A. C. Weideman.
The Exponentially Convergent Trapezoidal Rule.
SIAM Review 56 (2014), 385--458.
\url{https://people.maths.ox.ac.uk/trefethen/publication/PDF/2014_149.pdf}.

\bibitem{ALT}
K. M. R. Audenaert. On the Araki--Lieb--Thirring inequality.
\url{https://arxiv.org/abs/math/0701129}.

\bibitem{HP}
F. Hansen, G. K. Pedersen. Jensen's Operator Inequality.
Bull. London Math. Soc. 35 (2003), 553--564.
\url{https://arxiv.org/abs/math/0204049}.

\bibitem{Companion}
Exact small-exponent constants for fractional matrix maps.
Companion working manuscript, November 2026.
\href{sharp-fractional-holder.pdf}{Local manuscript}.
\end{thebibliography}
\end{document}
